\chapter{Python Bot 如何在线维护同一个状态机}

\section{本章要解决的困惑}

fast backtest 读 \code{decision_frame_v1}。Python Bot 在线运行时，输入不是 parquet，而是 Runner 发来的 public/private event stream。但它维护的是同一类状态：

\begin{lstlisting}
quote -> bid/ask/mid/spread/frames
trade -> rolling TFI
fill/private -> closed outcome memory R5
policy -> order intent
\end{lstlisting}

所以 Python Bot 不是另一套策略。它是同一个简单状态机的在线版本。

\section{online\_features.py}

代码入口：

\begin{lstlisting}
systems/ccusdt_replay_exchange/strategies/python/ccusdt_tfi_core_idle01/online_features.py
\end{lstlisting}

核心对象：

\begin{lstlisting}
RollingTrade
ClosedOutcomeState
OnlineFeatureState
\end{lstlisting}

\code{OnlineFeatureState.add_trade} 的逻辑和第一章一样：

\begin{lstlisting}
is_buy = side == "buy"
signed = qty if is_buy else -qty
self.signed_qty += signed
self.abs_qty += abs(qty)
\end{lstlisting}

TFI：

\begin{lstlisting}
def tfi(self) -> float:
    if self.abs_qty <= EPS:
        return 0.0
    return self.signed_qty / self.abs_qty
\end{lstlisting}

\section{update\_quote：frames 和 spread}

quote 进来后，Bot 更新：

\begin{lstlisting}
spread_bps = (ask - bid) / mid * 10000
mid_change = mid - last_mid
if mid changed:
    frames_since_mid_change = 0
else:
    frames_since_mid_change += 1
\end{lstlisting}

这和第一章的 L2 例子一致。区别是 fast 线已经把它算成 \code{decision_frame}，Bot 线要边收 quote 边维护。

\section{ClosedOutcomeState：R5 的在线来源}

R5 只应该来自已经闭合的交易：

\begin{lstlisting}
positive_bps = sum(value for value in outcomes_bps if value > 0)
negative_bps = sum(value for value in outcomes_bps if value < 0)
r5 = positive_bps / abs(negative_bps)
\end{lstlisting}

Bot 不能读取未来 PnL，也不能读取 \code{date/} 目录。它只能从 private fills 和已闭合 lifecycle 里更新状态。

\warningline{Bot 不读 date/。如果 Bot 读了研究输出，回测就不是 runtime-safe。}

\section{strategy.py：stdin/stdout sparse Bot}

\code{strategy.py} 是 Runner sparse-python 模式下的 Bot 入口。它接收：

\begin{lstlisting}
session_start
market_quote
market_trade
market_l2_update
market_batch
account_snapshot
order_ack
order_reject
fill
session_end
\end{lstlisting}

输出只应该是：

\begin{lstlisting}
heartbeat
submit_order
cancel_order
\end{lstlisting}

稀疏策略不应该对每个 market event 输出 hold。它只在需要下单或保活时输出。

\section{Bot 和 fast 的差别}

\begin{longtable}{p{0.22\textwidth}p{0.34\textwidth}p{0.34\textwidth}}
\toprule
对象 & fast backtest & Python Bot \\
\midrule
输入 & decision frame parquet & public/private stream \\
状态 & ShadowFourCellPolicy + cached frames & OnlineFeatureState + ShadowFourCellPolicy \\
输出 & entries/exits tables & submit order intent \\
成交 & fast 近似 & Runner 决定 \\
\bottomrule
\end{longtable}

fast 更适合研究参数，Bot 更接近 runtime。

\section{在线例子一：三笔 trade 如何变成 TFI}

Runner public stream 连续发来：

\begin{lstlisting}
market_trade side=buy  qty=300
market_trade side=buy  qty=200
market_trade side=sell qty=100
\end{lstlisting}

Bot 的 \code{OnlineFeatureState} 不需要知道完整订单簿，也不需要未来价格。它只维护 rolling 窗口：

\begin{lstlisting}
signed_qty = +300 +200 -100 = 400
abs_qty = 300 + 200 + 100 = 600
tfi = 400 / 600 = 0.667
\end{lstlisting}

这时当前旧 entry 不触发。若第三笔也是 buy：

\begin{lstlisting}
signed_qty = 600
abs_qty = 600
tfi = 1.0
\end{lstlisting}

才可能进入 policy trigger。在线和 fast 的区别只是数据来源不同，数学定义相同。

\section{在线例子二：quote 如何维护 frames}

Runner 发来 quote：

\begin{longtable}{p{0.18\textwidth}p{0.20\textwidth}p{0.20\textwidth}p{0.22\textwidth}}
\toprule
event & bid & ask & mid \\
\midrule
1 & 0.1120 & 0.1124 & 0.1122 \\
2 & 0.1120 & 0.1124 & 0.1122 \\
3 & 0.1120 & 0.1124 & 0.1122 \\
4 & 0.1121 & 0.1125 & 0.1123 \\
\bottomrule
\end{longtable}

Bot 维护：

\begin{lstlisting}
event 1: last_mid empty -> frames = 0
event 2: mid unchanged -> frames = 1
event 3: mid unchanged -> frames = 2
event 4: mid changed   -> frames = 0
\end{lstlisting}

这和 \code{decision_frame_v1} 的 \code{frames_since_mid_change} 是同一概念。Bot 不是从 parquet 把 frames 读出来，而是在线维护。

\section{在线例子三：fill 后如何形成 R5}

Bot 产生 order intent 后，成交事实由 Runner private stream 返回。只有在生命周期结束后，Bot 才能把结果放进 closed outcome memory。

\begin{lstlisting}
entry mid = 0.1122
close mid = 0.1125
direction = long
pnl_bps = log(0.1125 / 0.1122) * 10000 = 26.7
\end{lstlisting}

这个 \(26.7\) bps 才能成为后续 R5 的一个 \(g_i\)。如果 position 还没到 fixed60，或者 fill 尚未完成，它不能进入 R5。

\section{从 signal 到 submit order}

Bot 的 policy 如果给出 signal，后面仍要经过 admission 和 capacity：

\begin{lstlisting}
signal = ShadowFourCellPolicy.on_quote(...)
admission = EntrySpreadAdmissionGate.decide(signal)
capacity = CapacityLedger.decide(signal)
if admission.accepted and capacity.actual_exposure > 0:
    submit_order(...)
\end{lstlisting}

这段伪代码就是 fast runner 和 Bot 的共同骨架。差别在于：fast 写 \code{entries.parquet}，Bot 写 order intent 给 Runner。

\section{Bot 绝不能做的事}

当前 runtime safety 可以用一句话概括：Bot 只能看“到当前为止 Runner 发给它的东西”。所以它不能：

\begin{itemize}[leftmargin=2em]
  \item 读取 \code{date/} 里的未来 market 文件；
  \item 读取 \code{entries.parquet} 或 \code{exits.parquet}；
  \item 读取 scored entries、MFE、MAE、future return label；
  \item 为了调试把 future label 塞进 feature snapshot；
  \item 从 Monitor 接 order decision。
\end{itemize}

\warningline{Bot 只能产生策略意图。成交事实永远由 Runner 决定。}

\section{手动检查点}

读 Bot 代码时，按顺序看：

\begin{lstlisting}
online_features.py: add_trade, update_quote, tfi, ClosedOutcomeState
shadow_policy.py: on_quote / on_decision_frame
strategy.py: message loop and submit_order
execution_runtime.py: admission and capacity helpers
\end{lstlisting}

\checkpoint{Python Bot 不神秘。它只是把第一章的状态机放到事件流里在线维护。}
